\paragraph{} No contexto de otimização e logística, as instâncias de \textit{Solomon} constituem um conjunto de problemas de teste (\textit{benchmarks}) propostos por Marius Solomon em 1987 para o VRPTW. Essas instâncias são fundamentais para aferir a eficiência de algoritmos, pois representam diferentes desafios geográficos e temporais. As nomenclaturas \textbf{C}, \textbf{R} e \textbf{RC} referem-se à disposição espacial dos clientes:

\begin{itemize}
    \item \textbf{Tipo C (\textit{Clustered}):} Os clientes estão organizados em agrupamentos (\textit{clusters}). O algoritmo deve ser eficiente em atender aos clientes de um grupo denso antes de se deslocar para o próximo. Geralmente, este perfil é considerado o de mais fácil resolução para a obtenção de rotas ótimas.
    \begin{figure}[H]
        \centering
        \parbox[htb]{13.0cm}
        {
            \centering
            \includegraphics[scale=0.5]{C101_map.png}
            \caption[\small{Mapa da instância C101}]{\label{C101_map} \small{Distribuição geográfica dos clientes na instância \textit{Clustered} C101}}
        }
    \end{figure}
    \vspace{-0.5cm}
    
    \item \textbf{Tipo R (\textit{Random}):} Os clientes estão distribuídos de forma aleatória e uniforme no plano de coordenadas. A ausência de padrões óbvios exige que o algoritmo avalie criteriosamente as distâncias e as janelas de tempo para cada tomada de decisão.
    \begin{figure}[H]
        \centering
        \parbox[htb]{13.0cm}
        {
            \centering
            \includegraphics[scale=0.5]{R101_map.png}
            \caption[\small{Mapa da instância R101}]{\label{R101_map} \small{Distribuição geográfica dos clientes na instância \textit{Random} R101}}
        }
    \end{figure}
    \vspace{-0.5cm}
    
    \item \textbf{Tipo RC (\textit{Random-Clustered}):} Trata-se de um modelo híbrido que combina regiões de distribuição aleatória com agrupamentos definidos. Este cenário aproxima-se da realidade urbana complexa, onde centros comerciais densos convivem com áreas de entregas residenciais esparsas.
    \begin{figure}[H]
        \centering
        \parbox[htb]{13.0cm}
        {
            \centering
            \includegraphics[scale=0.5]{RC101_map.png}
            \caption[\small{Mapa da instância RC101}]{\label{RC101_map} \small{Distribuição geográfica dos clientes na instância \textit{Random-Clustered} RC101}}
        }
    \end{figure}
    \vspace{-0.5cm}
\end{itemize}



\paragraph{} Adicionalmente, as séries \textbf{1} e \textbf{2} definem o rigor das restrições operacionais, como a capacidade de carga dos veículos e a amplitude das janelas de tempo:

\begin{itemize}
    \item \textbf{Série 1:} Apresenta restrições severas, com veículos de baixa capacidade de carga e janelas de tempo estreitas. Tais limitações resultam em rotas curtas e exigem uma frota maior para cobrir todos os pontos. Como exemplo, a instância \textit{RC101} possui uma janela de tempo média de apenas $32,08$.
    \item \textbf{Série 2:} Oferece um cenário mais flexível, com veículos de grande capacidade e janelas de tempo expandidas. Um único veículo consegue atender a uma quantidade maior de clientes, transferindo o desafio principal para a minimização da distância total percorrida. A instância \textit{RC201}, utilizada nos testes deste trabalho, possui uma janela média significativamente superior, de $128,32$.
\end{itemize}